% =========================================================================
% ปฏิบัติการที่ 3: การเขียนโปรแกรมตรวจจับพิกัดเชิงพื้นที่และท่าทางมือ
% =========================================================================

\chapter{การเขียนโปรแกรมตรวจจับพิกัดเชิงพื้นที่และท่าทางมือ}
\label{ch:lab03}

\noindent{\large\bfseries\color{cyberblue} การคำนวณระยะทางแบบยูคลิดเพื่อตรวจจับการจีบนิ้วและการกรองสัญญาณรบกวน}

\vspace{0.4cm}

\begin{labobjectivebox}{การเขียนโปรแกรมตรวจจับพิกัดเชิงพื้นที่และท่าทางมือ}
\textbf{วัตถุประสงค์การเรียนรู้}
\begin{enumerate}[leftmargin=1.8cm, itemsep=2pt]
    \item เข้าใจการแปลงพิกัดภาพ 2 มิติสู่เวกเตอร์เชิงพื้นที่ 3 มิติ
    \item สามารถประยุกต์สมการระยะทางแบบยูคลิดเพื่อตรวจจับท่าทางจีบนิ้ว (Pinch Gesture)
    \item สามารถออกแบบตัวกรองแบบเอ็กซ์โพเนนเชียล (Exponential Moving Average) เพื่อลดการสั่นไหวของพิกัด
\end{enumerate}

\vspace{4pt}
\textbf{เครื่องมือและอุปกรณ์จำเป็น} \enspace MediaPipe Hands, โมดูลคณิตศาสตร์เวกเตอร์เชิงเส้น (Vector3)
\end{labobjectivebox}

\section{หลักการและทฤษฎีพื้นฐาน}

ในการสร้างระบบควบคุมไร้สัมผัส (Touchless Interaction) ท่าทางการจีบนิ้ว (Pinch Gesture) ระหว่างปลายนิ้วหัวแม่มือ (Landmark 4) และปลายนิ้วชี้ (Landmark 8) ถือเป็นท่าทางมาตรฐานเทียบเท่ากับการคลิกเมาส์ในระบบปฏิบัติการทั่วไป การตรวจจับท่าทางนี้ทำได้โดยการคำนวณระยะทางแบบยูคลิด 3 มิติ (3D Euclidean Distance) ระหว่างจุดทั้งสอง:
\begin{equation}
d_{\text{pinch}} = \sqrt{(x_8 - x_4)^2 + (y_8 - y_4)^2 + (z_8 - z_4)^2}
\end{equation}
เมื่อค่า $d_{\text{pinch}}$ มีค่าน้อยกว่าเกณฑ์ที่กำหนด (Threshold เช่น $d < 0.035$ ในระบบพิกัดสัมพัทธ์) ระบบจะส่งสัญญาณเหตุการณ์การคลิกหรือหยิบจับวัตถุ

อย่างไรก็ตาม พิกัดที่ได้จากระบบวิทัศน์คอมพิวเตอร์มักมีสัญญาณรบกวนความถี่สูง (High-frequency Jitter) ซึ่งเกิดจากความผันผวนของแสงและการประมวลผลระดับพิกเซล การแก้ปัญหานี้จำเป็นต้องใช้ตัวกรองปรับเรียบ (Smoothing Filter) เช่น Exponential Moving Average (EMA):
\begin{equation}
\mathbf{P}_{\text{filtered}}^{(t)} = \alpha \mathbf{P}_{\text{raw}}^{(t)} + (1 - \alpha) \mathbf{P}_{\text{filtered}}^{(t-1)}
\end{equation}
โดยที่ $\alpha \in (0, 1]$ คือค่าน้ำหนักปรับเรียบ (Smoothing Factor) หากกำหนด $\alpha$ ต่ำ สัญญาณจะนิ่งเรียบแต่มีความหน่วง (Latency) เพิ่มขึ้น

\begin{conceptbox}{สาระสำคัญของปฏิบัติการที่ 3}
การเข้าใจโครงสร้างทางคณิตศาสตร์และการประมวลผลเชิงพื้นที่ ช่วยให้นักศึกษาสามารถออกแบบระบบเสมือนจริงที่ตอบสนองต่อผู้ใช้งานได้อย่างเป็นธรรมชาติและแม่นยำสูง ทั้งยังสามารถต่อยอดสู่การพัฒนาแอปพลิเคชันทางวิทยาศาสตร์และอุตสาหกรรมได้อย่างมีประสิทธิภาพ
\end{conceptbox}

\section{ขั้นตอนการปฏิบัติการและการทดลอง}

ให้นักศึกษาดำเนินการสร้างไฟล์โครงการและเขียนคำสั่งตามลำดับขั้นตอนต่อไปนี้ โดยตรวจสอบความถูกต้องของไลบรารีที่เรียกใช้งานและเส้นทางไฟล์

\begin{codebox}{โครงสร้างโค้ดหลักของปฏิบัติการที่ 3}
\begin{lstlisting}[language=HTML]
// ฟังก์ชันคำนวณระยะทางแบบยูคลิด 3 มิติ
function calculateDistance(p1, p2) {
    const dx = p1.x - p2.x;
    const dy = p1.y - p2.y;
    const dz = p1.z - p2.z;
    return Math.sqrt(dx * dx + dy * dy + dz * dz);
}

// ตัวแปรเก็บพิกัดปรับเรียบก่อนหน้า
let smoothedIndex = { x: 0, y: 0, z: 0 };
const ALPHA = 0.35; // ค่าน้ำหนักตัวกรอง EMA
const PINCH_THRESHOLD = 0.035;

function processGestures(landmarks) {
    const thumbTip = landmarks[4];
    const indexTip = landmarks[8];

    // ปรับเรียบพิกัดนิ้วชี้ด้วยสมการ EMA
    smoothedIndex.x = ALPHA * indexTip.x + (1 - ALPHA) * smoothedIndex.x;
    smoothedIndex.y = ALPHA * indexTip.y + (1 - ALPHA) * smoothedIndex.y;
    smoothedIndex.z = ALPHA * indexTip.z + (1 - ALPHA) * smoothedIndex.z;

    const pinchDist = calculateDistance(thumbTip, indexTip);
    const isPinching = pinchDist < PINCH_THRESHOLD;

    return {
        cursor: smoothedIndex,
        isPinching: isPinching,
        distance: pinchDist
    };
}
\end{lstlisting}
\end{codebox}

\section{การทดสอบระบบและการบันทึกผล}

เมื่อรันโปรแกรมบนเซิร์ฟเวอร์จำลอง (Local Development Server) ให้ทำการทดสอบการทำงานตามเกณฑ์ประเมินในตารางต่อไปนี้ พร้อมบันทึกผลการสังเกตการณ์ลงในรายงานผลการทดลอง

\begin{table}[htbp]
\centering
\small
\begin{tabularx}{\textwidth}{c X c Y}
\toprule
\textbf{ลำดับ} & \textbf{รายการทดสอบและพฤติกรรมที่คาดหวัง} & \textbf{เกณฑ์ผ่าน} & \textbf{ผลการทดสอบ} \\
\midrule
1 & การเริ่มต้นระบบและการเข้าถึงสิทธิ์ฮาร์ดแวร์ & ภายใน 3 วินาที & ผ่าน / ไม่ผ่าน \\
2 & ความเสถียรของการตรวจจับและการคำนวณพิกัด & อัตราความแม่นยำ $> 90\%$ & ผ่าน / ไม่ผ่าน \\
3 & อัตราการแสดงผลกราฟิกต่อเนื่อง (Framerate) & ไม่ต่ำกว่า 55 FPS & ผ่าน / ไม่ผ่าน \\
4 & การตอบสนองต่อปฏิสัมพันธ์เชิงพื้นที่ & ความหน่วง $< 40$ ms & ผ่าน / ไม่ผ่าน \\
\bottomrule
\end{tabularx}
\caption{ตารางประเมินผลการทำงานของระบบในปฏิบัติการที่ 3}
\label{tab:eval_lab03}
\end{table}

\section{ภารกิจท้าทายและการต่อยอด}

\begin{activitybox}{กิจกรรมส่งเสริมทักษะขั้นสูง}
ให้นักศึกษาพัฒนาท่าทางเพิ่มเติมอีก 2 ท่าทาง ได้แก่ ท่าทางกำมือ (Fist Detection) โดยตรวจสอบระยะห่างระหว่างปลายนิ้วทั้งสี่กับจุดกึ่งกลางฝ่ามือ และท่าทางกางมือ (Open Palm) เพื่อใช้เป็นคำสั่งรีเซ็ตระบบ
\end{activitybox}

\section{คำถามท้ายการทดลอง}

ให้นักศึกษาตอบคำถามเชิงวิเคราะห์ต่อไปนี้ลงในสมุดบันทึกปฏิบัติการ

\begin{enumerate}[leftmargin=1.8cm, itemsep=6pt]
    \item เหตุใดการตรวจจับท่าทางจีบนิ้วจึงต้องคำนึงถึงพิกัดแกน $z$ ร่วมด้วย แทนที่จะคำนวณเฉพาะระยะทางบนระนาบ $x-y$
    \item จงวิเคราะห์ข้อดีและข้อเสียของการตั้งค่าพารามิเตอร์ $\alpha$ ในตัวกรอง EMA ที่มีค่าสูง ($0.8$) เทียบกับค่าต่ำ ($0.15$)
    \item หากระยะห่างระหว่างมือกับกล้องเปลี่ยนไป จะส่งผลกระทบต่อเกณฑ์การตัดสิน $d_{\text{pinch}}$ หรือไม่ อย่างไร
\end{enumerate}

